%3. Foreign aid: majority support for increased foreign aid if it is well spent
%============================================================================================================
%Reprendre/paraphraser/étendre et compléter (notamment avec mon papier NHB) la section A.1.3 de mon papier NHB
%============================================================================================================

\section{Foreign aid: majority support for increased foreign aid if it is well spent}

Public attitudes toward foreign aid in high-income countries have been studied for several decades. This literature shows that citizens commonly overestimate the amount their country spends on foreign aid, even though majorities generally support maintaining or increasing aid, particularly when it is perceived as effective. In the United States, 83\% of respondents supported a multilateral initiative to halve world hunger \citep{PIPA2001}. Across 20 countries, an average of 81\% agreed that developed countries ``have a moral responsibility to work to reduce hunger and severe poverty in poor countries'' \citep{PIPA2008}. Moreover, in each of seven OECD countries surveyed, at least 75\% of respondents were willing to contribute financially to a program designed to halve world hunger, for an estimated cost of, for example, \$50 per year for an American respondent.

%============================================================================================================
%- underestimation: Kaufman, Nair, Fabre et al NHB (Figs 8, S20, S23),
%============================================================================================================


\citet{kaufmann_foreign_2019} document that respondents in each of the 24 countries they study overestimate foreign aid, by a factor of seven on average. In most countries, the level of aid people would like their government to provide exceeds the level they believe is currently provided.\footnote{\citet{kaufmann_foreign_2012} provide the most reliable estimates of desired aid because, as noted by \citet{hudson_mile_2012}, earlier studies generally failed to account for misperceptions of actual aid levels.} Respondents in the top income quintile support aid levels that are 0.13 percentage points of GNI lower than those preferred by the bottom 40\%. Combining a theoretical framework with data on lobbying and observed aid allocations, while controlling for citizens' preferences in terms of desired aid, the authors conclude that actual aid falls short of desired aid because wealthier individuals exert disproportionate political influence toward their economic interests. The United States stands out as an exception. Desired aid is similar to the average across other countries, at around 3\% of GNI, but respondents believe that aid is already about twice as large as the level they would ideally support. Perceptions of current aid alone explain more than 15\% of the variation in desired aid. From one country to another, the ratio of median desired aid to actual aid ranges almost from parity to a ratio of more than 9 to 1, revealing what the authors describe as a substantial ``democratic aid deficit'' in several donor countries.

The United States is an important exception in the cross-country evidence. In the data analyzed by \citet{kaufmann_foreign_2019}, Americans preferred foreign aid equivalent to 3.19\% of GNI, close to the cross-country average, but believed that their country was already spending 7.52\%, compared with an actual level of only 0.18\%. This particularly large misperception is consistent with earlier U.S. evidence. \citet{gilens_political_2001} estimates that, among respondents with high levels of general political knowledge, informing them that foreign aid accounted for less than 1\% of federal spending reduced the predicted share favoring aid cuts from 58.6\% to 42.0\%. Likewise, \citet{nair_misperceptions_2018} attribute weak support for aid among US-Americans to widespread misperceptions of their rank in the global income distribution: respondents underestimate their rank by 27 percentiles on average while they overestimate the global median income by a factor of ten. Such misperceptions may explain why, in the 2000-2004 waves of the GSS, more than 60\% of Americans believed that the government was spending too much on foreign aid \citep{okten_preferences_2007}.

Using a survey experiments in the United States, the United Kingdom, France, Germany and Spain, \citet{fabre_majority_2025} show that, on average across these countries,
respondents greatly overestimate current foreign aid (e.g., 4.7\% against an actual 0.4\% of public spending in the U.S.), but still prefer substantially higher aid budgets than those currently observed, allocating on average 1.8\% after receiving information on the true amount and 4.0\% without it \citet[Figure 8]{fabre_majority_2025_si}.

Respondents were also asked how much they believed their country spent on foreign aid. The median respondent overestimated the actual amount. In the United States, the median estimate ranged from 1.8\% to 2.6\%, compared with an actual level of 0.4\%. In the pooled Western European sample, respondents believed that foreign aid represented 2.9\% of public spending, compared with an actual level of 1.1\%. Their mean preferred allocation was 3.9\% without information and 2.7\% after information was provided \citep[Figure S20]{fabre_majority_2025_si}.

%============================================================================================================
% - support Fabre et al NHB (Figs 6, 7, S21-S26)
%============================================================================================================

Respondents were then asked about their attitudes toward their country's evolution of foreign aid. 60\% of Americans are in favor for an increase, 43\% of them wants conditions attached to this, while the corresponding figures for the pooled Western European sample are 64\% and 39\% to 51\% \citep[Figure 6]{fabre_majority_2025_si}. In addition, respondents were asked about the government spending they would allocate to foreign aid, with a random subgroup first informed of the actual amount. Without this information, the median preferred interval was 1.8-2.6\% in the United States, France, and Germany, compared with 1.1-1.7\% in Spain and the United Kingdom. Informing respondents reduced the U.S. median to 0.3-0.5\%, while the French and German medians remained at 1.8-2.6\% \citep[Figure S20-S21]{fabre_majority_2025_si}. Respondents, for about 47\% in U.S., 59\% in Europe and even 75\% in France ``Preferred foreign aid is higher than current'' \citep[Figure S23]{fabre_majority_2025_si}. Respondents who preferred higher foreign aid were asked how the increase should be financed. ``Higher taxes on the wealthiest'' were by far the most frequently selected option, chosen by 68\% of Americans and 64\% of Europeans, and by as many as 82\% in Spain and 85\% in the United Kingdom. By contrast, relatively few, between 1\% and 8\% respondents favored cuts in education \citep[Figure S24]{fabre_majority_2025_si}. Respondents who preferred lower aid were asked how the budgetary savings should be used instead. Healthcare services spending was generally prioritized: 57\% of Europeans favored it while it represents 40\% of Americans \citep[Figure S25]{fabre_majority_2025_si}.


%============================================================================================================
% - conditions: primarily no corruption
%============================================================================================================

Respondents who are in favor of an increase in foreign aid only under certain conditions were asked: ``Under which conditions should foreign aid be increased?''. \citet{fabre_majority_2025} find that the most important condition is ensuring that ``the aid reaches people in need and money is not diverted'', supported by 77\% of Europeans and 68\% of Americans. A second widely supported condition is ``that recipient countries comply with climate targets and human rights'' (72\% in Europe and 61\% in the United States). Around half of respondents also require that ``other high-income countries also increase their foreign aid'' (51\% in Europe and 45\% in the U.S.), whereas conditioning aid on ``increased taxes on millionaires'' receives more limited support (38\% and 36\%, respectively) \citep[Figure 7]{fabre_majority_2025}.

Reviewing the literature, \citet{hudson_mile_2012} conclude that support for foreign aid is largely motivated by an intrinsic concern for poverty, while also documenting considerable skepticism about aid effectiveness. According to \citet{dfid_aid_2009} and \citet{PIPA2001}, 47\% of British respondents believe that aid is wasted, primarily because of corruption, while Americans estimate that less than one-quarter of aid reaches people in need and that more than half is captured by corrupt government officials. These beliefs do not, however, rule out a utilitarian motivation for supporting aid: even if only a fraction of the funds reaches the intended recipients, the expected welfare gains from transferring resources to substantially poorer populations may remain large \citep{kopczuk_limitations_2005}.

Consistent with \citet{henson_public_2010}, \citet{bauhr_corruption_2013} find that perceived corruption in recipient countries lowers support for aid. However, this relationship is mitigated by what they describe as the aid-corruption paradox: countries where corruption is major are often those with the highest need for assistance. \citet{bodenstein_who_2017} further show that right-leaning Europeans and respondents who perceive high levels of corruption are more likely to support conditioning EU development assistance on recipient countries complying with standards of democracy, human rights, and good governance.

Respondents who preferred reducing or maintaining foreign aid were asked: ``Why should foreign aid should not be increased?'' The most common reason is ``Charity begins at home: there is already a lot to do to support the [country] people in need'', selected by 63\% of Europeans and Americans who refuse to increase aid. The second most frequent answer is ``Aid is not effective as most of it is diverted'', chosen by slightly half of these respondents %53\% of Europeans and 40\% of Americans 
\citep[Figures 7, S25]{fabre_majority_2025}.%LB: or: \citep[Figure S25]{fabre_majority_2025_si} for the citation 


%============================================================================================================
% - determinants: higher support on the left and among rich, greater effect of ideology among the rich (Cho et al 24)
%============================================================================================================
%LB: la dernière fois, on avait parlé de (Cho, 2024) dont les résultats en intro sont les suivants: « My empirical findings indicate that financial circumstances often outweigh this normative impact of ideology, especially when respondents are explicitly directed to take their own financial situation into account when indicating the level of support they favor. My results also show that left-leaning citizens are not always more supportive of sharing resources with poor people abroad than are right-leaning voters. » (Cho, 2024, p. 542)
% --> donc je pense que je vais juste utiliser la Figure 1 page 10. Et peut être les données de l'Eurobaromètre comme tu l'avais dit (sauf si je les mobilise au début de la section). % AF oui très bien

Using the 2002 Gallup survey and the 2006 World Values Survey, \citet{bayram_aiding_2017} and \citet{paxton_individual_2012}, consistent with the U.S. evidence of \citet{heinrich_voters_2018}, identify trust, left-wing political orientation, political interest, and female gender as the main individual-level correlates of support for higher aid.

% AF: ici tu montres que parmi les gens de gauche, les riches sont plus en faveur que les pauvres. Est-ce qu'on peut montrer plutôt la réciproque, à savoir l'effet de l'idéologie au sein des riches par rapport aux pauvres ? %LB: est-ce que tu pensais à quelque chose comme ça ? 
Using Eurobarometer data from 15 OECD donor countries, \citet{cho_ideology_2024} shows that the association between ideology and the intensity of moral support for foreign aid is stronger among socioeconomically advantaged respondents than among disadvantaged respondents. At the far-left end of the ideological scale, the predicted probability of answering ``totally agree'' that tackling poverty abroad is a moral obligation is approximately 63\% among advantaged respondents, compared with 44\% among disadvantaged respondents. Conversely, the probability of answering ``tend to agree'' is around 31\% among advantaged respondents and 42\% among disadvantaged respondents. Thus, among the advantaged left, the highest level of support clearly predominates, whereas among the disadvantaged left, ``totally agree'' and ``tend to agree'' are nearly equally prevalent. As respondents move to the right, ``tend to agree'' becomes the most prevalent response among advantaged respondents, while its predicted probability remains relatively stable across the ideological scale among disadvantaged respondents. Ideology therefore more strongly differentiates the intensity of support among advantaged respondents, whereas disadvantaged respondents tend to express more moderate support across ideological positions \citep[Figure 1]{cho_ideology_2024}.

%LB: ajout ici de l'Eurobaromètre
Eurobarometer surveys reveal broad acceptance of foreign aid across the European Union. In 2015, in the 28 EU Member States, approximately 64\% supported honouring the EU's commitment to increase aid or increasing it beyond the promised amount \citep{european_commission_2015_eurobarometer}. 
In 2019, 73\% thought that the EU and its Member States should either spend more on financial assistance to partner countries (30\%) or maintain spending at its current level (43\%). Of the remaining respondents, 17\% favored lower spending and 10\% answered ``Don't know'' \citep{european_commission_2019_development}.
In 2020, a survey across the 27 Member States found that 91\% of respondents considered EU funding for humanitarian aid important \citep{european_commission_2021_humanitarian}.
In 2022, in all 27 Member States, 80\% stated poverty reduction in developing countries should be a priority of the EU, and 67\% considered it should be a priority for their national government \citep{european_commission_2022_development}. % AF  Am I right to shift to 'should'? %LB: oui (p. 21: "should be a priority for their national government")
